% ============================================================================
%  Section I — Introduction  (IOP MST; passive voice / third person)
% ============================================================================

\section{Introduction}
\label{sec:intro}

Induction machines account for a dominant share of the electrical
energy consumed by industry and form the mechanical backbone of countless
processes; their unexpected failure causes costly downtime and, in critical
applications, safety hazards. Continuous condition monitoring (CM) and early
fault diagnosis are therefore central to modern maintenance
practice~\cite{gangsar2020signal, hamani2025advancements}. Among the available
modalities, vibration analysis is the most established, because the mechanical
and electrical faults of interest---bearing defects, rotor asymmetries, and
supply imbalances---imprint characteristic signatures on the machine's vibration
response~\cite{yang2021fault, zhang2020deep}.

The prevailing paradigm relies on industrial-grade piezoelectric accelerometers
sampled at tens to hundreds of kilohertz, which resolve the high-frequency
bearing defect frequencies exploited by spectral and envelope analysis. This is
reflected in the public benchmarks that anchor the field, acquired at
$8$--$200$~kHz~\cite{lee2024multidomain, sehri2024ottawa, thuan2023hust,
huang2018bearing}. Although accurate, such instrumentation is costly, requires
intrusive installation, and scales poorly to the vast population of small,
low-value machines that individually do not justify a dedicated monitoring channel
yet collectively dominate a plant's machine inventory. Closing this coverage gap
motivates a shift toward inexpensive, wireless, and readily deployable sensing.

Consumer micro-electro-mechanical-systems (MEMS) inertial sensors, epitomised by
those embedded in smartphones, offer such a platform at negligible marginal
cost~\cite{grouios2022accelerometers}. Machine fault detection has been demonstrated
from smartphone- and MEMS-acquired vibration~\cite{naha2017mobile,
ertargin2023deep, ertargin2025automated}, and dedicated low-cost MEMS nodes now
integrate on-board processing with wireless telemetry for machine
diagnostics~\cite{pedotti2020lowcost, jakobsen2024lowcost}. These sensors,
however, operate at sampling rates one to three orders of magnitude below
industrial practice---$100$~Hz in the present case---so that, by the Nyquist
criterion~\cite{nyquist1924certain, shannon1948mathematical}, the high-frequency
signatures underpinning classical bearing diagnosis are not observable. A
fundamental and largely unresolved question follows: what fault information is
genuinely recoverable from such a constrained sensor, and does the apparent
performance transfer to a real deployment?

The second half of that question is easily overlooked. Learning-based fault
diagnosis routinely reports near-perfect accuracy, yet a growing body of evidence
attributes much of this success to \emph{data leakage}---information shared
between training and test partitions that inflates measured accuracy without
improving generalisation~\cite{kapoor2023leakage}. In vibration diagnosis, leakage
arises when overlapping analysis windows from a single continuous recording are
divided across partitions, or when recording-specific characteristics are learned
in place of fault features; enforcing component- or recording-wise separation has
been shown to reduce reported accuracy substantially~\cite{wheat2024impact,
hendriks2022towards}. This scrutiny has so far been directed at industrial-grade,
high-rate data. For low-rate smartphone datasets---in which each operating
condition is captured in one continuous acquisition and heavy window overlap is
common---the risk is arguably greater, yet random-split evaluation remains the
norm, including in the baseline accompanying the dataset examined
here~\cite{ertargin2026smartphone}.

A third consideration is deployment. The Industrial Internet of Things envisions
dense networks of inexpensive nodes that stream, or locally distil, diagnostic
information over constrained wireless links~\cite{yousuf2024iot}, and recent
TinyML implementations have shown that lightweight classifiers can run on
microcontrollers at millisecond latency and sub-millijoule
energy~\cite{gao2025edge, asutkar2023tinyml}. Such work establishes that edge
diagnosis is feasible, but it generally presupposes industrial-bandwidth sensing
and does not characterise the coupled trade-off between diagnostic accuracy and
the sampling, communication, and computational cost that governs an ultra-low-rate
node. Whether, for a smartphone-class sensor, the low sampling rate is a liability
to be tolerated or an asset to be exploited has not been quantified.

This paper approaches these gaps as a measurement problem. Using a public
smartphone-acquired vibration dataset of an induction machine recorded across six
health states, three supply frequencies, and two load conditions at
$100$~Hz~\cite{ertargin2026smartphone, ertargin2026dataset}, and rather than
proposing a new classifier, the study foregrounds measurement rigour: the
consumer sensor is characterised as a measurement instrument, a protocol is
developed for obtaining unbiased and uncertainty-quantified estimates of its
diagnostic performance, and the accuracy--resource trade-off that governs an IoT
deployment is mapped. Two measurement questions organise the work---how much
fault information the sensor can convey, and how that performance must be measured
if it is to transfer to deployment. The principal contributions of this paper are
fourfold. First, it introduces a \emph{performance-measurement methodology} for
windowed low-rate vibration condition monitoring---comprising leakage-free
recording-wise and cross-condition (cross-speed and cross-load) protocols with
paired significance testing and expanded-uncertainty reporting---and applies it to
quantify the measurement bias introduced by the prevailing random-split practice,
which is shown to overstate accuracy by approximately $41$ to $47$ percentage
points. Second, it provides a \emph{measurement characterisation} of the
smartphone MEMS channel that establishes, with uncertainty, its effective
resolution, signal floor, and usable band, and reveals that the on-device gravity
compensation preserves the vibration power but reshapes its low-frequency
structure, leaving the raw channels the more informative and hence preferable for
condition monitoring. Third, it develops an
\emph{accuracy--resource characterisation} that identifies a resource-frugal
operating point which, under leakage-free evaluation, classifies the machine's
health state significantly more accurately than the full-rate configuration, and
it translates
this trade-off, through a communication link budget and an edge-model footprint,
into deployable operating points for Bluetooth Low Energy, Zigbee, NB-IoT, and
LoRaWAN nodes. Finally, it releases a \emph{reproducible and uncertainty-aware
baseline} for low-cost IoT-based induction-machine condition monitoring, with all
code and configurations made available.

The remainder of this paper is organised as follows.
Section~\ref{sec:related} reviews related work on vibration-based condition
monitoring, low-cost MEMS sensing, IoT and edge diagnosis, and evaluation
methodology. Section~\ref{sec:dataset} describes the dataset and measurement
setup, and Section~\ref{sec:charact} characterises the sensor and its signals.
Section~\ref{sec:method} details the processing chain and the evaluation
protocols. Section~\ref{sec:results} reports the results,
Section~\ref{sec:discussion} discusses their implications and limitations, and
Section~\ref{sec:conclusion} concludes.
